\documentclass[12pt]{article}
\usepackage{graphicx}

\usepackage{hyperref}  %allows for clickable reference
\hypersetup{colorlinks = true,citecolor = blue, linkcolor = blue, urlcolor = blue}
\usepackage{subcaption}

\usepackage{amsmath}
\usepackage{pifont}
\usepackage{wrapfig}
\usepackage{multicol}
\setlength{\columnsep}{1.5cm}
\setlength{\columnseprule}{1pt}
 
\usepackage{tikz}
\usetikzlibrary{calc}


\usepackage{fancyhdr}

\lhead{ }
\rhead{ }
\chead{Assignment-6}
 
\renewcommand{\headrulewidth}{0.4pt} %thickness of Line
\rfoot{Page \thepage}
\fancyfoot[CO,RE]{Vivek Kumar}
\fancyfoot[LO,CE]{04-Mar-2022}
\renewcommand{\footrulewidth}{0.2pt}  %thickness of line

 
\begin{document}
\pagestyle{empty}

\begin{tikzpicture}
[remember picture, overlay]
\draw [line width=2pt]($(current page.north west)+(1in,-1in)$) rectangle ($(current page.south east)+(-1in, 1in)$);

\end{tikzpicture}


   \begin{center}
       \vspace{1cm}
	   \Large
       \textbf{National Institute Of Technology, Raipur } 
       \vspace{1.5cm}
     
       \includegraphics[scale=0.1]{nit-raipur-logo.png}\\
       \vspace{0.8cm}
       \Huge
       Assignment-6\\
       \vspace{0.8cm}
      \textbf{ "Two page write-up on 5 Solutions to Covid-19 provided by Biomedical Engineers."}
      
	\vfill      
      
   \begin{multicols}{2} 
   \begin{flushleft}
       \large
       \textbf{Submitted by:}\\
       Name: Vivek Kumar\\
       Roll no: 21111071\\
       Branch : Biomedical\\
       Semester : $1^{st}$ \\
       NIT Raipur,Chhattisgarh\\
       \columnbreak
       \textbf{Under the Supervision of:}\\
       Mr.Saurabh Gupta
       Department of Biomedical Engineering
       NIT Raipur, Chhattisgarh
    \end{flushleft}
    \end{multicols}    
            
   \end{center}

\clearpage
\pagestyle{fancy}
\tableofcontents
\clearpage
\section{Introduction}
From the Thermometer measuring Fever to the ventilators helping keep patients alive, Biomedical engineers have a big role to play in humanity's efforts to fight COVID-19. The opportunity to help society is what engineering is all about, said \href{https://www.eng.famu.fsu.edu/cbe/people/locke}{Bruce Locke}, a distinguished research professor of chemical and biomedical engineering at the FAMU-FSU College of Engineering.
Many people working in engineering have responded to the ongoing crisis by adapting their existing skills and equipment to help fight COVID-19. The role of a Biomedical Engineer includes designing biomedical equipment and devices to aid the recovery or improve the health of individuals.
 
\section{Solutions to Covid-19 provided by Biomedical Engineers}
\begin{itemize}
\item \textbf{Rapid POC and home-based testing/diagnostics} : Home based testing kits for corona virus plays a vital role in determining the presence of infection when in the first time. This can roughly confirm that you have corona infection or not. This is a quick testing process which can give almost immediate result to check your infection.
Home testing kits for COVID-19 which is based on saliva can be much more effective because that is how official testing is taking place. It may be possible that person going through that test may become negative in test, so it is recommended to test it once again on a gap of 3-4 days.

\item \textbf{Portable Emergency Venrtilator} : Many COVID-19 patients experience Acute Respiratory Distress Syndrome (ARDS), a condition that can be treated with mechanical ventilation. In response to the need for mechanical ventilators, designed and tested an emergency ventilator (EV) that can control a patient's peak inspiratory pressure (PIP) and breathing rate, while keeping a positive end expiratory pressure (PEEP). Engineers From Various Countries and Institutions came up with Different Prototypes in Order to Develope Low Cost, Efficient and Portable EV.

\item \textbf{Digital health platforms and models that integrate data, Management Tools} : Big data and artificial intelligence (AI) have helped 
facilitate COVID-19 preparedness and the tracking of 
people, and so the spread of infection, in several 
countries. Tools such as migration maps, which use 
Mobile Phones Network, GPS, and Social 
Media to collect real-time data on the location of people.

\item \textbf{Nanotechnology Biosensor Being Adapted for Rapid COVID-19 Testing} : Srikanth Singamaneni, professor of mechanical engineering and materials science, and his team have developed a rapid, highly sensitive and accurate biosensor based on an ultrabright fluorescent nanoprobe, which has the potential to be broadly deployed, Called plasmonic-fluor. The ultrabright fluorescent nanoprobe can also help in resource-limited conditions because it requires fewer complex instruments to read the results.

\href{https://teknosenin.com/newly-developed-nanotechnology-biosensor-being-adapted-for-rapid-covid-19-testing/}{Singamaneni hypothesizes their plasmonic-fluor-based biosensor} will be 100 times more sensitive compared with the conventional SARS-CoV-2 antibody detection method. Increased sensitivity would allow clinicians and researchers to more easily find positive cases and lessen the chance of false negatives.

\item \textbf{Technologies for protecting healthcare workers, first responders, and caregivers} : The protection of health care workers is vital in continuing patient care in health care systems that are currently challenged by the pandemic, but also important in ensuring they do not spread the virus. so, access to personal protective equipment (PPE) for health workers is a key concern. Therefore PPE Kit, Medical Masks and Gloves are Manufactured in a Large Scale across the Globe to Help in Fighting with Pandemic.

\end{itemize}
 

\section{Conclusion}
This pandemic has helped to highlight some of the unseen professions that help to make our health service work, and show the positive impact that Biomedical Engineering in particular can have on people's lives. During This Covid-19 event Biomedical Engineering has suddenly become the most important engineering discipline in the world.

\end{document}